\section{Preliminary}
\label{sec:preliminary}

\subsection{Tasks and hypothesis spaces}

We adopt the formulation of inductive bias learning introduced by
\citet{baxter2000model}. Let $\mathcal X$ and $\mathcal Y$ denote the input
and output spaces, and let $\ell:\mathcal Y\times\mathcal Y\to[0,\infty)$
be a loss function. A supervised learning task is represented by a
probability distribution $P$ on $\mathcal X\times\mathcal Y$.
A hypothesis is a function $h:\mathcal X\to\mathcal Y$, whose population
risk is
\begin{equation}
  R_P(h)
  :=\mathbb E_{(x,y)\sim P}\bigl[\ell(h(x),y)\bigr].
  \label{eq:population-risk}
\end{equation}
We assume that the expectations appearing below are well defined and finite.

For a single task, a task learner receives a sample
$S=\{(x_j,y_j)\}_{j=1}^{m}$ drawn
independently from $P$. Its empirical risk is
\begin{equation}
  \widehat R_S(h)
  :=\frac{1}{m}\sum_{j=1}^{m}\ell(h(x_j),y_j).
  \label{eq:empirical-risk}
\end{equation}
A task learner $A_{\mathcal H}$ maps this sample to a hypothesis
$A_{\mathcal H}(S)\in\mathcal H$ within a given hypothesis space
$\mathcal H$. The choice of $\mathcal H$
encodes an inductive bias: restricting the available functions can
facilitate generalization from limited data, provided that the space
retains sufficiently good solutions.

\subsection{Task environment and bias learner}

Let $Q$ be a probability distribution over a collection $\mathcal P$ of
task distributions. Source tasks $P_1,\ldots,P_n$ are drawn independently
from $Q$, and each supplies a sample $S_i$ of independent observations.
Given a family of hypothesis spaces $\mathbb H$, a bias learner $B$ uses
these samples to select a shared hypothesis space:
\begin{equation}
  \widehat{\mathcal H}
  :=B(S_1,\ldots,S_n)\in\mathbb H,
  \label{eq:bias-learner}
\end{equation}
Task learners use the selected space to learn their own hypotheses on
new tasks drawn independently from $Q$. Following Baxter, the suitability
of a space for this environment is measured by
\begin{equation}
  R_Q(\mathcal H)
  :=\mathbb E_{P\sim Q}
    \left[\inf_{h\in\mathcal H}R_P(h)\right].
  \label{eq:environment-risk}
\end{equation}
The task-specific infimum allows different hypotheses to solve different
tasks within the same shared space.

\subsection{Connectivity as a structural inductive bias}

In our setting, hypothesis spaces are indexed by connectivity masks.
Fix a network architecture with $E$ candidate connections in the layer
whose structure is to be learned. A binary mask
\begin{equation}
  M\in\mathcal M\subseteq\{0,1\}^{E}
  \label{eq:admissible-masks}
\end{equation}
specifies which connections are available. The admissible set
$\mathcal M$ may impose a fixed connection budget or additional
connectivity constraints.

Let $f_{w,M}$ denote a predictor with mask $M$ and numerical parameters
$w$. Masked connections are set to zero, while $w$ includes the active
weights and the remaining trainable parameters. Each mask defines a
hypothesis space
\begin{equation}
  \mathcal H_M
  :=\{f_{w,M}:w\in\Theta_M\},
  \qquad
  \mathbb H_{\mathcal M}
  :=\{\mathcal H_M:M\in\mathcal M\},
  \label{eq:masked-hypothesis-spaces}
\end{equation}
where $\Theta_M$ is the admissible parameter set.

Learning a structural bias therefore amounts to selecting a mask whose
induced space is suitable for the task environment. The corresponding
population objective is
\begin{equation}
  \inf_{M\in\mathcal M}R_Q(\mathcal H_M)
  =\inf_{M\in\mathcal M}\mathbb E_{P\sim Q}
    \left[\inf_{w\in\Theta_M}R_P(f_{w,M})\right].
  \label{eq:structural-bias-objective}
\end{equation}
In this setting, the bias learner $B$ selects a mask $\widehat M$ and
thereby returns the space $\widehat{\mathcal H}=\mathcal H_{\widehat M}$.
The selected mask is fixed across tasks, while a task learner trains
numerical parameters from scratch for each new task. Connectivity constrains the
support of the weights; equality constraints between active weights
require an additional parameter-sharing mechanism.

\subsection{Learning new tasks from finite samples}

The environment risk measures the best solutions available within a
hypothesis space. To describe learning from limited data, we also specify
how a task learner uses that space.

For a fixed mask $M$, let $A_M$ denote the task learner operating within
$\mathcal H_M$. It returns a predictor $A_M(S)\in\mathcal H_M$
using sample $S$. Its expected risk after observing $m$ labeled examples is
\begin{equation}
  \mathcal R_m(M)
  :=\mathbb E_{P\sim Q}\,
    \mathbb E_{S\sim P^m}
      \left[R_P\bigl(A_M(S)\bigr)\right].
  \label{eq:finite-sample-transfer-risk}
\end{equation}
For a mask $\widehat M$ selected by the bias learner $B$, this quantity
is evaluated conditional on the source data and the selected mask.
The new task and its sample are independent of this selection.

Our objective is to learn a fixed structure that improves
$\mathcal R_m$ at limited sample sizes, or reaches a prescribed risk
level with fewer examples than comparison structures under a common
training protocol. A low value of $R_Q(\mathcal H_M)$ establishes that
good solutions are available. Whether the training procedure can obtain
them efficiently from finite samples is a separate question.
